\scanpage{226}

\mspage{100.}

donc $C\cap W=C\cap W^{\mathrm{univ}}$, \deleted{\unsure{qui}} \unsure{qui} est stable par cochgt. de base puisque $C$ et $W^{\mathrm{univ}}$ le sont. Quant à 5°) c'est une tautologie, \deleted{compte} tenu de 4°.

8. \uline{Application aux catégories de modèles} $(\widehat A,W_A)$.

Soit $A$ dans $\mathrm{Cat}$, et considérons dans $\widehat A$ le localiseur $W_A$, image inverse par

\[
\left\{\begin{aligned}i_A:\quad&\widehat A\longrightarrow\mathrm{Cat}\\&F\longmapsto A/F\end{aligned}\right.\tag*{(8.1.)}
\]

d'un localiseur fondamental donné $W$, satisfaisant les conditions

\[
W\bigl(\underset{\text{sat.}}{1},\underset{\text{loc.}}{2}, \underset{\substack{\text{\unsure{dépl.}}\\\text{fin.}}}{3}, \underset{\substack{\text{\unsure{conn}}\\\Pi_0}}{5}, \underset{\varinjlim}{6}, \underset{\substack{\text{carrés}\\\text{\unsure{coc.} en}\\\text{``M.V.''}}}{7\,\mathrm{bis}\deleted{\unsure{3}}}, \underset{\text{acc.}}{8}.\bigr)\tag*{(8.2.)}
\]

\margin{où $W(8)$ est l'accessibilité de $W$ (donc $W$ est L-accessible s'il satisfait aussi $W(\unsure{6|5})$).}

\deleted{\unsure{Comme} $\widehat A$ \unsure{est}} Pour pouvoir \unsure{à présent|maintenant} appliquer les développements sur les parties accessibles d'une grosse catégorie, \deleted{il faut} du \S\,5, il faut vérifier les conditions a) b) c) d) e) \deleted{\unsure{\S}} du \S\,5, tant \deleted{dans} $\mathrm{Cat}$, que (surtout) dans $\widehat A$.

a) Stabilité par \deleted{\unsure{petites}} \uline{$\lim$} filtrantes, OK.

b) Filtration cardinale. Bien connue pour $\widehat A=\underline{\mathrm{Hom}}\,(A,\mathrm{Ens})$, à partir de celle de $\mathrm{Ens}$ par les cardinaux [cf SGA 4 I 9.2\unsure{4|.4}, p. 172]. Pour $\mathrm{Cat}$, on définit une filtration cardinale

\[
(\mathrm{Cat}^{\Pi})_{\Pi\geq\Pi_0},\quad\text{où }\Pi_0=\text{Cardinal des ens. dénombrables}\tag*{(8.3)}
\]
% archive sheet number: 225
